\chapter{Goniometrie in de rechthoekige driehoek}\label{ch:g9:trig}

In een \omterm{def:g6:shapes:triangles}{rechthoekige driehoek} ligt de \emph{vorm} van de driehoek vast zodra je één
scherpe hoek kent --- alleen zijn grootte kan nog veranderen. De verhoudingen van
zijn zijden hangen dus alleen van de hoek af: die verhoudingen zijn de \omterm{def:g9:trig:ratios}{cosinus}, de
\omterm{def:g9:trig:ratios}{sinus} en de \omterm{def:g9:trig:ratios}{tangens}. Samen met de stelling van Pythagoras laten ze je elke zijde en
elke hoek van een \omterm{def:g6:shapes:triangles}{rechthoekige driehoek} uit heel weinig gegevens berekenen.

\section{De stelling van Pythagoras, opnieuw}

\begin{theorem}[Pythagoras]\label{thm:g9:trig:pythagoras}
In een driehoek $ABC$ die \omterm{def:g6:shapes:triangles}{rechthoekig} is in $A$, is het kwadraat van de schuine
zijde de \omterm{def:g1:addition:def}{som} van de kwadraten van de twee rechthoekszijden:
\[
BC^2 = AB^2 + AC^2 .
\]
Omgekeerd: geldt $BC^2 = AB^2 + AC^2$ in een driehoek, dan is de driehoek
\omterm{def:g6:shapes:triangles}{rechthoekig} in $A$.
\end{theorem}
\admitted

\begin{example}\label{ex:g9:trig:pythagoras}
Een \omterm{def:g6:shapes:triangles}{rechthoekige driehoek} heeft rechthoekszijden $AB = 5$ en $AC = 12$. Dan is
\[
BC^2 = 5^2 + 12^2 = 25 + 144 = 169,
\qquad\text{dus}\quad BC = \sqrt{169} = 13 .
\]
Om een \emph{rechthoekszijde} te vinden: is $BC = 10$ en $AB = 6$, dan is
$AC^2 = BC^2 - AB^2 = 100 - 36 = 64$, dus $AC = 8$.
\end{example}

\section{Cosinus, sinus, tangens}

In een \omterm{def:g6:shapes:triangles}{rechthoekige driehoek} hebben de drie zijden, ten opzichte van een scherpe
hoek $\theta$, elk een naam: de \emph{schuine zijde} (tegenover de \omterm{def:g3:shapes:rightangle}{rechte hoek}, de
langste zijde), de zijde \emph{overstaand} aan $\theta$, en de zijde
\emph{aanliggend} aan $\theta$ (de rechthoekszijde die aan $\theta$ raakt).

\begin{definition}[Goniometrische verhoudingen]\label{def:g9:trig:ratios}
Voor een scherpe hoek $\theta$ van een \omterm{def:g6:shapes:triangles}{rechthoekige driehoek}:
\[
\cos\theta = \frac{\text{aanliggend}}{\text{schuine zijde}},
\qquad
\sin\theta = \frac{\text{overstaand}}{\text{schuine zijde}},
\qquad
\tan\theta = \frac{\text{overstaand}}{\text{aanliggend}} .
\]
Deze verhoudingen\index{cosinus}\index{sinus}\index{tangens} hangen alleen van
$\theta$ af, niet van de grootte van de driehoek (alle \omterm{def:g6:shapes:triangles}{rechthoekige driehoeken} met
dezelfde scherpe hoek zijn schalingen van elkaar, volgens Thales).
\end{definition}

\begin{omfigure}
\begin{tikzpicture}[scale=1.15]
  \coordinate (A) at (0,0);
  \coordinate (B) at (4.2,0);
  \coordinate (C) at (4.2,2.2);
  \draw[thick] (A) -- (B) -- (C) -- cycle;
  \draw[thin] (B) ++(-0.28,0) -- ++(0,0.28) -- ++(0.28,0);
  \draw[omThm, ->] (0.9,0) arc (0:27.7:0.9);
  \node[omThm, font=\small] at (1.25,0.21) {$\theta$};
  \node[below, font=\small, omDef] at (2.1,0) {aanliggend};
  \node[right, font=\small, omMeth] at (4.2,1.1) {overstaand};
  \node[above left=-2pt, font=\small, omProp] at (2.2,1.25) {schuine zijde};
  \node[below left, font=\small] at (A) {$B$};
  \node[below right, font=\small] at (B) {$A$};
  \node[above right, font=\small] at (C) {$C$};
\end{tikzpicture}

{\small De drie zijden gezien vanuit de hoek $\theta = \widehat B$: de schuine
zijde $[BC]$, de aanliggende zijde $[BA]$, de overstaande zijde $[AC]$.}
\end{omfigure}

\begin{proposition}[Eerste eigenschappen]\label{prop:g9:trig:properties}
Voor elke scherpe hoek $\theta$:
\begin{enumerate}
  \item $0 < \cos\theta < 1$ en $0 < \sin\theta < 1$ (een rechthoekszijde is
        korter dan de schuine zijde);
  \item $\cos^2\theta + \sin^2\theta = 1$;
  \item $\tan\theta = \dfrac{\sin\theta}{\cos\theta}$.
\end{enumerate}
\end{proposition}

\begin{proof}
Schrijf $a$ voor de aanliggende zijde, $o$ voor de overstaande zijde en $h$ voor de
schuine zijde van een \omterm{def:g6:shapes:triangles}{rechthoekige driehoek} met hoek $\theta$.

1. $0 < a < h$ en $0 < o < h$, dus liggen de twee verhoudingen strikt tussen $0$ en
$1$.

2. Volgens Pythagoras is $a^2 + o^2 = h^2$. Deel alles door $h^2$:
\[
\cos^2\theta + \sin^2\theta
= \frac{a^2}{h^2} + \frac{o^2}{h^2}
= \frac{a^2 + o^2}{h^2} = 1 .
\]

3. $\dfrac{\sin\theta}{\cos\theta} = \dfrac{o/h}{a/h} = \dfrac oa
= \tan\theta$.
\end{proof}

\begin{method}[Een zijde vinden]\label{met:g9:trig:side}
Om een onbekende zijde te berekenen wanneer één zijde en één scherpe hoek bekend
zijn:
\begin{enumerate}
  \item markeer de bekende hoek en benoem de drie zijden (schuine zijde,
        overstaand, aanliggend) \emph{vanuit die hoek};
  \item kies de verhouding (cos, sin of tan) waarin de bekende zijde en de
        gevraagde zijde voorkomen;
  \item schrijf de \omterm{def:g8:equations:def}{vergelijking} op en los ze op;
  \item controle op het gezond verstand: de schuine zijde moet er als langste uit
        komen.
\end{enumerate}
\end{method}

\begin{example}\label{ex:g9:trig:side}
In een \omterm{def:g6:shapes:triangles}{rechthoekige driehoek} meet de schuine zijde $8$ en is één hoek $35^\circ$;
bereken de zijde die daar tegenover ligt. De verhouding waarin de overstaande zijde
en de schuine zijde voorkomen, is de \omterm{def:g9:trig:ratios}{sinus}:
\[
\sin 35^\circ = \frac{\text{overstaand}}{8}
\quad\Longrightarrow\quad
\text{overstaand} = 8 \sin 35^\circ \approx 8 \times 0.574 \approx 4.6 .
\]
Voor de aanliggende zijde gebruik je de \omterm{def:g9:trig:ratios}{cosinus}:
$8\cos 35^\circ \approx 8 \times 0.819 \approx 6.6$. Controle met Pythagoras:
$4.6^2 + 6.6^2 \approx 21 + 43 \approx 64 = 8^2$.
\end{example}

\begin{method}[Een hoek vinden]\label{met:g9:trig:angle}
Zijn twee zijden bekend, bereken dan de verhouding die ze vormen, bepaal of die de
\omterm{def:g9:trig:ratios}{cosinus}, de \omterm{def:g9:trig:ratios}{sinus} of de \omterm{def:g9:trig:ratios}{tangens} van de onbekende hoek is, en pas de omgekeerde
functie van de rekenmachine ($\cos^{-1}$, $\sin^{-1}$ of $\tan^{-1}$) toe op de
verhouding.
\end{method}

\begin{example}\label{ex:g9:trig:angle}
Een ladder van $5$ m staat tegen een muur, met zijn voet $1.4$ m van de muur. De
hoek $\theta$ tussen de ladder en de grond voldoet aan
\[
\cos\theta = \frac{1.4}{5} = 0.28,
\qquad\text{dus}\quad
\theta = \cos^{-1}(0.28) \approx 73.7^\circ .
\]
\end{example}

\begin{omfigure}
\begin{tikzpicture}[scale=0.85]
  \draw[thick] (0,0) -- (1.4,0);
  \draw[thick] (1.4,0) -- (1.4,4.0);
  \draw[omDef, very thick] (0,0) -- (1.4,4.0);
  \draw[thin] (1.4,0) ++(-0.25,0) -- ++(0,0.25) -- ++(0.25,0);
  \draw[omThm, ->] (0.75,0) arc (0:70.7:0.75);
  \node[omThm, font=\small] at (1.15,0.55) {$\theta$};
  \node[below, font=\small] at (0.7,0) {$1.4$ m};
  \node[omDef, font=\small, left] at (0.62,2.1) {$5$ m};
  \draw[gray, thin] (-0.7,0) -- (2.3,0);
\end{tikzpicture}

{\small Het ladderprobleem: de aanliggende zijde ($1.4$ m) en de schuine zijde
($5$ m) zijn bekend, dus geeft de \omterm{def:g9:trig:ratios}{cosinus} de hoek.}
\end{omfigure}

\begin{example}[Bijzondere hoeken]\label{ex:g9:trig:special}
Twee driehoeken geven exacte waarden. De \omterm{def:g3:division:half}{helft} van een vierkant met zijde 1 (een
rechthoekige \omterm{def:g6:shapes:triangles}{gelijkbenige driehoek}) heeft schuine zijde $\sqrt2$ en hoeken van
$45^\circ$:
\[
\cos 45^\circ = \sin 45^\circ = \frac{1}{\sqrt2} = \frac{\sqrt2}{2},
\qquad \tan 45^\circ = 1 .
\]
De \omterm{def:g3:division:half}{helft} van een \omterm{def:g6:shapes:triangles}{gelijkzijdige driehoek} met zijde 1 geeft de hoeken $30^\circ$ en
$60^\circ$:
\[
\cos 60^\circ = \sin 30^\circ = \frac12,
\qquad
\cos 30^\circ = \sin 60^\circ = \frac{\sqrt3}{2}.
\]
\end{example}

\section{Oefeningen}

\begin{exercise}[$\star$]\label{exo:g9:trig:1}
Een \omterm{def:g6:shapes:triangles}{rechthoekige driehoek} heeft rechthoekszijden $9$ en $12$. Bereken zijn schuine
zijde. Een andere heeft schuine zijde $17$ en één rechthoekszijde $8$: bereken de
andere rechthoekszijde.
\end{exercise}

\begin{exercise}[$\star$]\label{exo:g9:trig:2}
Een driehoek heeft zijden $7$, $24$ en $25$. Is hij \omterm{def:g6:shapes:triangles}{rechthoekig}? Dezelfde vraag
voor zijden $5$, $6$ en $8$.
\end{exercise}

\begin{exercise}[$\star$]\label{exo:g9:trig:3}
In een driehoek $DEF$ die \omterm{def:g6:shapes:triangles}{rechthoekig} is in $D$, met de hoek in $E$ aangeduid als
$\theta$: welke zijde is de schuine zijde? Welke zijde ligt tegenover $\theta$?
Schrijf $\cos\theta$, $\sin\theta$ en $\tan\theta$ als verhoudingen van de zijden
$DE$, $DF$, $EF$.
\end{exercise}

\begin{exercise}[$\star$]\label{exo:g9:trig:4}
In een \omterm{def:g6:shapes:triangles}{rechthoekige driehoek} meet de schuine zijde $10$ en is één scherpe hoek
$28^\circ$. Bereken de twee rechthoekszijden (afgerond op een tiende;
$\cos 28^\circ \approx 0.883$, $\sin 28^\circ \approx 0.469$).
\end{exercise}

\begin{exercise}[$\star$]\label{exo:g9:trig:5}
In een \omterm{def:g6:shapes:triangles}{rechthoekige driehoek} meet de rechthoekszijde die aan de hoek $\theta$ ligt
$6$, en de overstaande rechthoekszijde $4.5$. Bereken $\tan\theta$, en daarna
$\theta$ (afgerond op de graad; $\tan^{-1}(0.75) \approx 36.9^\circ$).
\end{exercise}

\begin{exercise}[$\star\star$]\label{exo:g9:trig:6}
Een drone vliegt op een hoogte van $120$ m. Een waarnemer op de grond ziet hem
onder een hoogtehoek van $32^\circ$. Op welke horizontale afstand van de waarnemer
bevindt de drone zich ($\tan 32^\circ \approx 0.625$)?
\end{exercise}

\begin{exercise}[$\star\star$]\label{exo:g9:trig:7}
Een oprit moet $1.2$ m stijgen met een hoek van ten hoogste $5^\circ$ met de
horizontale. Welke minimale lengte langs de helling moet de oprit hebben
($\sin 5^\circ \approx 0.0872$)? Rond naar boven af op een tiende meter.
\end{exercise}

\begin{exercise}[$\star\star$]\label{exo:g9:trig:8}
Met \cref{prop:g9:trig:properties}: een scherpe hoek $\theta$ voldoet aan
$\cos\theta = 0.6$. Bereken $\sin\theta$ zonder $\theta$ te zoeken, en daarna
$\tan\theta$.
\end{exercise}

\begin{exercise}[$\star\star\star$]\label{exo:g9:trig:9}
Verantwoord de exacte waarden van \cref{ex:g9:trig:special}: bereken in de
rechthoekige \omterm{def:g6:shapes:triangles}{gelijkbenige driehoek} met rechthoekszijden $1$ de schuine zijde;
bereken in de \omterm{def:g6:shapes:triangles}{gelijkzijdige driehoek} met zijde $1$ de hoogte, en leid de \omterm{def:g9:trig:ratios}{cosinus} en
de \omterm{def:g9:trig:ratios}{sinus} van $30^\circ$ en $60^\circ$ af.
\end{exercise}

\section{Opgave: meten wat je niet kunt bereiken}

\begin{problem}[{Weekendopgave --- de methode met twee standplaatsen van de
landmeter voor onbereikbare hoogten, hellingspercentages, de afstand tot de
horizon, en de wedloop van de tangens naar het oneindige}]
\label{pb:g9:trig:1}
Geen enkel meetlint reikt tot de top van een berg, tot de spits van een toren aan
de andere kant van een rivier, of tot de horizon op zee --- maar een hoekmeter en
de drie verhoudingen van dit hoofdstuk wel. Het middelpunt van deze opgave is de
klassieke \emph{methode met twee standplaatsen} van de landmeters, die een hoogte
berekent zonder ooit haar voet te benaderen; eromheen: verkeersborden, ladders,
trappen, de kromming van de Aarde, en een eerste glimp van een functie die naar het
oneindige ontploft.

\medskip
\noindent\textbf{Deel I --- Eén standplaats.}
(Waarden: $\tan 35^\circ \approx 0.700$, $\sin 35^\circ \approx 0.574$,
$\cos 35^\circ \approx 0.819$, $\tan 1^\circ \approx 0.0175$.)
\begin{enumerate}
  \item Vanuit een punt op $60$ m van de voet van een vuurtoren wordt de top gezien
        onder een hoogtehoek van $35^\circ$ (gemeten vanaf de horizontale, op
        ooghoogte). Hoe hoog ligt de top boven ooghoogte?
  \item Een verkeersbord kondigt een stijging van $12\,\%$ aan: de weg stijgt
        $12$ m per $100$ m horizontaal. Welke hoek maakt de weg met de horizontale?
        En welk stijgingspercentage zou een weg van $45^\circ$ hebben?
  \item De twee scherpe hoeken van een \omterm{def:g6:shapes:triangles}{rechthoekige driehoek} zijn complementair, en
        de overstaande zijde van de ene is de aanliggende zijde van de andere. Leid
        de ``co''-gelijkheden af:
        \[
        \sin(90^\circ - \theta) = \cos\theta,
        \qquad
        \cos(90^\circ - \theta) = \sin\theta .
        \]
  \item Een ladder van $8$ m staat onder $60^\circ$ met de grond. Geef met de
        exacte waarden van \cref{ex:g9:trig:special} de bereikte hoogte en de
        afstand van de voet tot de muur --- exact, en daarna op de centimeter.
  \item Ga $\cos^2\theta + \sin^2\theta = 1$ met getallen na voor
        $\theta = 35^\circ$ en exact voor $\theta = 60^\circ$. Waarom is deze
        gelijkheid (\cref{prop:g9:trig:properties}) een goede gewoonte om je werk
        met de rekenmachine te controleren?
\end{enumerate}

\medskip
\noindent\textbf{Deel II --- Twee standplaatsen.}
De top van een berg is zichtbaar, zijn voet onbereikbaar. Vanuit een punt $A$ is de
hoogtehoek van de top $30^\circ$; na $200$ m rechtdoor naar de berg te lopen, naar
$B$, wordt de hoogtehoek $40^\circ$. Schrijf $h$ voor de hoogte van de top boven
ooghoogte en $x$ voor de horizontale afstand van $B$ tot de verticale van de top.
($\tan 30^\circ \approx 0.577$, $\tan 40^\circ \approx 0.839$.)
\begin{enumerate}[resume]
  \item Druk $\tan 40^\circ$ en $\tan 30^\circ$ uit met $h$, $x$ en $x + 200$: twee
        \omterm{def:g8:equations:def}{vergelijkingen} voor twee onbekenden.
  \item Druk uit elke \omterm{def:g8:equations:def}{vergelijking} de horizontale afstand in $h$ uit, trek af, en
        leid de formule met twee standplaatsen af:
        \[
        h = \frac{200}
        {\ \dfrac{1}{\tan 30^\circ} - \dfrac{1}{\tan 40^\circ}\ } .
        \]
        Bereken $h$ op tien meter nauwkeurig.
  \item Bereken ook $x$, en toets je twee antwoorden aan de peiling van $30^\circ$
        vanuit $A$.
  \item In één zin: waarom is de methode met twee standplaatsen onmisbaar --- welke
        ene meting, die hier onmogelijk is, zou de methode met één standplaats uit
        vraag 1 hebben gevraagd?
  \item Een toren wordt gepeild onder $30^\circ$, en onder $60^\circ$ na $200$ m
        naar hem toe te lopen. Bereken de hoogte \emph{exact} met
        $\tan 30^\circ = \frac{1}{\sqrt3}$ en $\tan 60^\circ = \sqrt3$ --- het
        antwoord is een \omterm{def:g9:arith:divisor}{veelvoud} van $\sqrt3$ --- en daarna op de meter.
\end{enumerate}

\medskip
\noindent\textbf{Deel III --- Tot de horizon en tegen de muur op.}
\begin{enumerate}[resume]
  \item Hoe ver ligt de horizon? Staat je oog $h$ meter boven een bolvormige Aarde
        met \omterm{def:g3:shapes:shapes}{radius} $R$, dan scheert je blik langs de bol: de bliklijn, de \omterm{def:g3:shapes:shapes}{radius}
        naar het scherende punt en de \omterm{def:g3:shapes:shapes}{radius} naar je voeten vormen een \omterm{def:g6:shapes:triangles}{rechthoekige
        driehoek}. Toon met Pythagoras (\cref{thm:g9:trig:pythagoras}) aan dat de
        afstand $d$ tot de horizon voldoet aan $d^2 = 2Rh + h^2 \approx 2Rh$, en
        bereken $d$ voor ogen op $1.80$ m ($R = 6.37 \times 10^6$ m).
  \item Dezelfde vraag vanaf de top van een toren van $324$ m. (Vuistregel van
        zeelieden: de horizon in kilometer is ongeveer $3.6\sqrt{h}$ met $h$ in
        meter --- toets je twee antwoorden eraan.)
  \item Je duim, ongeveer $2$ cm breed, op armlengte gehouden, ongeveer $60$ cm van
        je oog: welke hoek bedekt hij? De volle Maan bedekt ongeveer $0.5^\circ$:
        welk deel van je duim verbergt de hele Maan?
  \item Comfortabele trappen hebben optreden van $17$ cm en aantreden van $29$ cm.
        Onder welke hoek klimmen ze? Vergelijk met de $30^\circ$ van
        \cref{ex:g9:trig:special}, en met een steile zoldertrap onder $60^\circ$:
        welke optrede zou \emph{die} nodig hebben bij een aantrede van $29$ cm?
  \item De wedloop van de \omterm{def:g9:trig:ratios}{tangens}: bereken (met de rekenmachine)
        $\tan 10^\circ$, $\tan 45^\circ$, $\tan 80^\circ$, $\tan 89^\circ$. Wat
        gebeurt er als de hoek $90^\circ$ nadert, en waarom (denk aan de
        aanliggende zijde)? De muur staat verticaal, de verhouding heeft geen
        waarde, en de grafiek schiet naar het oneindige --- je eerste
        \emph{asymptoot}, een wezen dat in het bovenbouwvolume nauwkeurig wordt
        bestudeerd.
\end{enumerate}
\end{problem}
